\chapter{La anticoncepción}\label{ch:g8:contraception}

Los capítulos anteriores cartografiaron un mecanismo de una potencia
notable: dos \omterm{def:g5:first-look-at-cells:cell}{células} se encuentran y, cuarenta semanas después, hay una
persona. Las personas somos la única \omterm{def:g5:biodiversity:species}{especie} que entiende su propio
mecanismo --- y entenderlo trae una elección a la que la rana no se
enfrenta nunca: si ponerlo en marcha, y cuándo. La
\emph{\omterm{prop:g8:contraception:principle}{anticoncepción}} es la ciencia aplicada de esa elección, y
funciona, en todos sus métodos, interrumpiendo un paso al que ya le
puedes poner nombre.

\section{El principio: nombra el paso, bloquea el paso}

\begin{proposition}[La anticoncepción se corresponde con el mecanismo]\label{prop:g8:contraception:principle}
El \omterm{def:g5:human-reproduction:pregnancy}{embarazo} exige una cadena sin roturas
(\cref{prop:g8:sexual-reproduction:core},
\cref{prop:g8:from-egg-to-birth:firstweek}): \omterm{def:g8:sexual-reproduction:gametes}{gametos} producidos ---
\omterm{def:g8:reproductive-systems:male}{espermatozoides} depositados --- el viaje a nado hasta el \omterm{def:g8:reproductive-systems:female}{oviducto} ---
una \omterm{def:g5:first-look-at-cells:cell}{célula} \omterm{def:g4:animal-reproduction:sexual}{femenina} soltada y encontrada --- \omterm{def:g5:flower-to-fruit:steps}{fecundación} ---
implantación. La \emph{anticoncepción}\index{anticoncepción} es
cualquier interrupción deliberada de esta cadena. Todos los métodos que
hay en la estantería de la farmacia son un ataque a un eslabón con
nombre, y clasificarlos por su eslabón es toda la teoría:
\begin{itemize}
  \item los métodos \emph{de barrera} impiden que se encuentren los
        \omterm{def:g8:sexual-reproduction:gametes}{gametos};
  \item los métodos \emph{hormonales} impiden que se suelte siquiera la
        \omterm{def:g5:first-look-at-cells:cell}{célula} \omterm{def:g4:animal-reproduction:sexual}{femenina};
  \item otros métodos hacen que falle el encuentro o la implantación.
\end{itemize}
\end{proposition}
\admitted

\section{Los métodos principales, por su eslabón}

\begin{definition}[Métodos de barrera: el preservativo]\label{def:g8:contraception:condom}
El \emph{preservativo}\index{preservativo} --- una funda fina que se
pone en el \omterm{def:g8:reproductive-systems:male}{pene} (hay una versión \omterm{def:g4:animal-reproduction:sexual}{femenina} que forra la vagina) ---
bloquea la cadena en el paso del depósito: el semen no entra nunca en
el cuerpo de la pareja. Está disponible en todas partes, hace falta
solo cuando hace falta --- y es único en la estantería por un segundo
servicio: es el único método que además bloquea el paso de las
\emph{infecciones} que se transmiten con el sexo, el mundo del
\cref{ch:g9:microbes-and-infection}. Esa doble tarea es la razón de que
el preservativo siga en la foto aunque se esté usando otro método.
\end{definition}

\begin{definition}[Métodos hormonales: la píldora]\label{def:g8:contraception:pill}
La \emph{píldora anticonceptiva}\index{píldora anticonceptiva} usa la
propia maquinaria de la \cref{def:g8:puberty:hormones}: tomada a
diario, sus \omterm{def:g8:puberty:hormones}{hormonas} mantienen el ciclo de los \omterm{def:g8:reproductive-systems:female}{ovarios} en el estado
suspendido (el acontecimiento 4 de la
\cref{prop:g8:reproductive-systems:cycle} sin ningún \omterm{def:g5:human-reproduction:pregnancy}{embarazo}) ---
\emph{sin \omterm{prop:g8:reproductive-systems:cycle}{ovulación}, sin \omterm{def:g5:first-look-at-cells:cell}{célula} \omterm{def:g4:animal-reproduction:sexual}{femenina}, no hay nada que fecundar}. La
receta un médico, que ajusta el método a la persona; es muy fiable si
se toma con regularidad --- y no ofrece absolutamente nada contra las
infecciones, de ahí la sociedad permanente con el \omterm{def:g8:contraception:condom}{preservativo}.
\end{definition}

\begin{example}[La estantería, clasificada por eslabón]\label{ex:g8:contraception:shelf}
Los demás métodos habituales, archivados por el eslabón que atacan: los
implantes y los parches hormonales --- el principio de la píldora con
mejor memoria (meses o años, sin toma diaria); los dispositivos
pequeños que un médico coloca en el \omterm{prop:g5:human-reproduction:plan}{útero} --- que hacen fallar el
encuentro y la implantación durante años seguidos; y la \emph{píldora
de urgencia}, una dosis hormonal de después del hecho, para tomar en
unos días, que funciona sobre todo retrasando una \omterm{prop:g8:reproductive-systems:cycle}{ovulación} que todavía
no ha ocurrido --- un freno de último recurso, disponible en la
farmacia, y ningún sustituto de un método usado \emph{antes}. Lo que no
hace ninguno de ellos: nada contra las infecciones.
\end{example}

\begin{remark}[Lo que no funciona]\label{rem:g8:contraception:notwork}
Los métodos del patio, medidos contra la cadena: ``contar los días
seguros'' se hunde con la variabilidad de la
\cref{rem:g8:reproductive-systems:living} --- la fecha de la \omterm{prop:g8:reproductive-systems:cycle}{ovulación}
no se puede saber por adelantado y los \omterm{def:g8:reproductive-systems:male}{espermatozoides} sobreviven días
esperándola; ``retirarse a tiempo'' se hunde con el hecho de que los
\omterm{def:g8:reproductive-systems:male}{espermatozoides} viajan antes del final. Los dos fallan justo donde
predice la teoría de los eslabones con nombre: no bloquean ningún
eslabón de manera fiable. La leyenda no es un método.
\end{remark}

\section{Elección, responsabilidad y el mapa de la ayuda}

\begin{proposition}[La elección es real y compartida]\label{prop:g8:contraception:choice}
Los hechos que impone el mecanismo:
\begin{itemize}
  \item el \omterm{def:g5:human-reproduction:pregnancy}{embarazo} necesita un solo acto, no la duración de una
        relación: la cadena no cuenta ocasiones;
  \item es una cadena \emph{compartida}: \omterm{def:g5:first-look-at-cells:cell}{células} de los dos miembros de
        la pareja, responsabilidad de los dos --- el \omterm{def:g8:contraception:condom}{preservativo} y la
        píldora no son cada uno asunto de un sexo, sino una
        conversación;
  \item y ningún método salvo la abstinencia es perfecto, aunque los
        buenos se acercan mucho si se usan exactamente como están
        diseñados --- en el ``bien usado'' es donde viven casi todos
        los fallos del mundo real.
\end{itemize}
\end{proposition}
\admitted

\begin{example}[El mapa de la ayuda]\label{ex:g8:contraception:help}
Las direcciones del \cref{ex:g8:puberty:questions}, ampliadas: las
enfermeras escolares y los médicos de familia responden preguntas sobre
\omterm{prop:g8:contraception:principle}{anticoncepción} con confidencialidad --- también a menores; los centros
de planificación familiar existen en todas las regiones justamente para
esto, gratis y sin juzgar; y los farmacéuticos manejan el reloj de la
píldora de urgencia. El sistema está hecho para que nadie tenga que
navegar la cadena a base de leyenda --- el mapa de la ayuda forma parte
del método.
\end{example}

\begin{method}[Evaluar cualquier método del que oigas hablar]\label{met:g8:contraception:evaluate}
Ante cualquier método que se presente como tal, pregunta por orden:
\begin{enumerate}
  \item ¿\emph{qué eslabón} de la cadena bloquea --- se lo puedes
        nombrar en el mapa?
  \item ¿con \emph{qué fiabilidad}, usado como lo usa la gente de
        verdad?
  \item ¿protege también contra las \emph{infecciones} o tiene que
        acompañarlo el \omterm{def:g8:contraception:condom}{preservativo}?
  \item ¿qué exige --- memoria diaria, receta, una colocación --- y le
        va eso a quien lo va a usar?
\end{enumerate}
Sin un eslabón que se pueda nombrar no hay método (la
\cref{rem:g8:contraception:notwork} falla en la pregunta 1).
\end{method}

\section{Ejercicios}

\begin{exercise}[$\star$]\label{exo:g8:contraception:1}
Enuncia la cadena que exige el \omterm{def:g5:human-reproduction:pregnancy}{embarazo} y define la \omterm{prop:g8:contraception:principle}{anticoncepción}
frente a ella.
\end{exercise}

\begin{exercise}[$\star$]\label{exo:g8:contraception:2}
¿Qué eslabón bloquea el \omterm{def:g8:contraception:condom}{preservativo}? ¿Y la píldora?
\end{exercise}

\begin{exercise}[$\star$]\label{exo:g8:contraception:3}
¿Qué segundo servicio da únicamente el \omterm{def:g8:contraception:condom}{preservativo}?
\end{exercise}

\begin{exercise}[$\star$]\label{exo:g8:contraception:4}
¿Cómo funciona sobre todo la píldora de urgencia y qué no es?
\end{exercise}

\begin{exercise}[$\star$]\label{exo:g8:contraception:5}
Nombra tres direcciones del mapa de la ayuda, con lo que ofrece cada
una.
\end{exercise}

\begin{exercise}[$\star$]\label{exo:g8:contraception:6}
¿Por qué no es fiable ``contar los días seguros''? Dos razones de la
\cref{rem:g8:contraception:notwork}.
\end{exercise}

\begin{exercise}[$\star\star$]\label{exo:g8:contraception:7}
Clasifica la estantería del \cref{ex:g8:contraception:shelf} por el
eslabón atacado, añadiendo el \omterm{def:g8:contraception:condom}{preservativo} y la píldora.
\end{exercise}

\begin{exercise}[$\star\star$]\label{exo:g8:contraception:8}
¿Por qué dicen tan a menudo los médicos ``píldora \emph{y}
\omterm{def:g8:contraception:condom}{preservativo}'' a las parejas jóvenes? Dos protecciones, una frase para
cada una.
\end{exercise}

\begin{exercise}[$\star\star$]\label{exo:g8:contraception:9}
Explica, con la \cref{prop:g8:reproductive-systems:cycle}, en qué
estado mantiene la píldora el ciclo y por qué eso evita el \omterm{def:g5:human-reproduction:pregnancy}{embarazo}.
\end{exercise}

\begin{exercise}[$\star\star$]\label{exo:g8:contraception:10}
``\omterm{def:g5:first-look-at-cells:cell}{Células} de los dos, asunto de los dos''. Defiende la frase frente a
la costumbre de dejarle la \omterm{prop:g8:contraception:principle}{anticoncepción} a un solo lado de la pareja.
\end{exercise}

\begin{exercise}[$\star\star$]\label{exo:g8:contraception:11}
Pásale el \cref{met:g8:contraception:evaluate} a: (a) el implante
hormonal; (b) ``retirarse a tiempo''.
\end{exercise}

\begin{exercise}[$\star\star\star$]\label{exo:g8:contraception:12}
``La \omterm{prop:g8:contraception:principle}{anticoncepción} es fisiología aplicada: cada método de la
estantería es un capítulo de este libro, apuntado''. Justifícalo con
tres métodos y sus tres capítulos.
\end{exercise}

\section{Problema: La consulta de planificación familiar}

\begin{problem}[{Problema de fin de semana --- una tarde de preguntas
en el centro de planificación}]\label{pb:g8:contraception:1}
La tarde (inventada) de un orientador de planificación familiar: cuatro
visitantes, cuatro situaciones y un mapa de la cadena en la pared.

\medskip\noindent\textbf{Parte I --- El mapa de la pared.}
\begin{enumerate}
  \item Dibuja el mapa de la pared: la cadena de los \omterm{def:g8:sexual-reproduction:gametes}{gametos} a la
        implantación, con los tres puntos de ataque de la
        \cref{prop:g8:contraception:principle} señalados.
  \item La visitante 1 pregunta ``¿cómo va a parar nada una pastilla al
        día?''. Responde con la maquinaria: ¿qué idioma habla la
        píldora y a quién?
  \item La visitante 1 insiste: ``¿y si se me olvida alguna?''. Localiza
        el peligro en el mapa --- ¿qué se reanuda cuando falla la señal
        hormonal?
  \item ¿Por qué sobrevive la recomendación del \omterm{def:g8:contraception:condom}{preservativo} del
        orientador a cualquier otra elección de método? (La doble
        tarea.)
\end{enumerate}

\medskip\noindent\textbf{Parte II --- Las situaciones.}
\begin{enumerate}[resume]
  \item La visitante 2, de diecisiete años y angustiada, creía que una
        aplicación de días seguros hacía innecesarios los demás
        métodos. Corrige la creencia con amabilidad: ¿qué dos hechos
        biológicos derrotan al calendario?
  \item El visitante 3 quiere algo ``sin memoria diaria''. Ofrécele los
        dos candidatos de la estantería y sus eslabones.
  \item La visitante 4 llega antes de un día desde que se rompió un
        \omterm{def:g8:contraception:condom}{preservativo}. Nombra la opción, su reloj, su mecanismo
        principal --- y el método permanente que el orientador va a
        comentar para después.
  \item Los cuatro se van con la misma frase final sobre las
        infecciones. Escríbela.
\end{enumerate}

\medskip\noindent\textbf{Parte III --- La reunión posterior.}
\begin{enumerate}[resume]
  \item La persona en prácticas pregunta por qué el centro enseña la
        cadena en vez de repartir métodos y ya está. Responde con la
        primera pregunta del \cref{met:g8:contraception:evaluate}.
  \item Clasifica los métodos de la tarde según lo que le exige cada
        uno a quien lo usa (la pregunta 4 del método).
  \item La persona en prácticas se pregunta por qué las visitas de los
        menores son confidenciales por diseño. Responde con la lógica
        del \cref{ex:g8:contraception:help}.
  \item Cierra la tarde en una frase: qué reparte en realidad el
        centro, más allá de las cajas.
\end{enumerate}
\end{problem}
